\documentclass[10pt,a4paper]{amsart}
\usepackage[margin=19mm]{geometry}
\usepackage{amsmath,amssymb,mathtools}
\usepackage[scr=rsfs,cal=euler]{mathalfa}
\usepackage{xfrac}
\usepackage[unicode,hidelinks,bookmarks=false]{hyperref}
\newcommand{\ie}{i.e.\ }
\newcommand{\cf}{cf.\ }
\newcommand{\resp}{respectively\ }
\newcommand{\Subsection}{\subsection}
\providecommand{\nobreakdash}{\nobreak\hbox{-}}
\setlength{\emergencystretch}{2em}
\multlinegap0pt
\usepackage{amssymb}
\usepackage{enumerate}
\usepackage[shortlabels]{enumitem}
\usepackage{bm}
\usepackage[normalem]{ulem}
\usepackage{xcolor}
\newtheorem{proposition}{Proposition}
\newcommand\aAd[3][E]{\sA_{#3}(#2|#1)}%conditional aggregation gdy uzywamy indeksów dólnych
\newcommand\bF{\mathbf{F}} %the set of functions
\newcommand\bFs{\bF_0^2} %the subset of FxF.
\newcommand\bMf{\mathbf{M}_{\mathrm{f}}}%the set of family of onotone measure
\newcommand{\bmu}{\boldsymbol{\mu}}
\newcommand\bmub[1][\bpeA]{\bmu_{#1}}
\newcommand\bpeA{\mathscr {A}_{\sbullet|\sbullet}}%family of CAO
\newcommand\cE{{\mathcal E}}
\renewcommand\ge{\geqslant}
\renewcommand\le{\leqslant}
\newcommand\sA{\mathsf{A}}
\newcommand\sbullet[1][.5]{\mathbin{\vcenter{\hbox{\scalebox{#1}{$\bullet$}}}}}
\title{Generalized level measure based on a~family \\of conditional aggregation operators}
\author{Michal Boczek, Ondrej Hutn\'ik, Marek Kaluszka, Miriam Kleinov\'a}
\date{Source: arXiv:2204.03069v1 (CC BY 4.0)}
\begin{document}
\maketitle

\noindent\emph{Source and licence.} Generalized level measure based on a~family \\of conditional aggregation operators. Version arXiv:2204.03069v1 (CC BY 4.0), distributed by arXiv under the Creative Commons Attribution 4.0 International licence, http://creativecommons.org/licenses/by/4.0/, as recorded in the arXiv licence metadata for this exact version and shown on the abstract page https://arxiv.org/abs/2204.03069v1. Full licence notice: \texttt{licenses/arxiv-2204.03069-CC-BY-4.0.txt}.\par\smallskip

\noindent\textit{Editorial scope.} Counted unit: the article's proposition pro:3.19 (source0.tex lines 725--731). The article's complete original proof of it is delivered with it (source0.tex lines 732--739, one \texttt{proof} environment of 8 lines). No mathematics is written, repaired or re-derived: the statement and the proof are the article's own lines, in the article's own notation, and no retained line is substituted or re-flowed.  No further source-local context is needed: every cross-reference of every retained line resolves inside the two delivered spans.  Nothing was omitted from any retained span, so the package carries no omission record.  No retained line contains a citation, so the wrapper carries no bibliography and no bibliography entry.  No retained line contains a figure environment, an image-inclusion command or drawing code, so no image is delivered and the package declares no asset dependency.  Editorial formatting only, in addition: an AMS \texttt{amsart} preamble that restates the article's own declarations and macros used by the retained lines, namely the environments \texttt{proposition} on separate counters, together with the standard packages those lines need.  The retained environments print in this document as Proposition 1 in order of appearance, whereas the article prints the same items with the numbering of its own full document.  The complete native source of the article as distributed in the arXiv e-print for this exact version is bundled unchanged as \texttt{sources/124241/source0.tex}.
\begin{proposition}\label{pro:3.19}
If $\bmu\in\bMf$ is nonincreasing, and  $\bpeA$ is a~$c$-quasi-subadditive on $\bF_0^2$ and nonincreasing pFCA, then
\begin{align}\label{eq:3.19}
\bmub(f+g,ct)\le \bmub(f,\lambda t)\vee \bmub(g,(1-\lambda)t)
\end{align}
for any $t\ge 0,$  any $\lambda\in (0,1),$  and  any $(f,g)\in\bFs.$
\end{proposition}

\begin{proof} Since $c\ge 1$ and pFCA $\bpeA$ is nonincreasing, we have $\cE_{ct}\subseteq \cE_t\subseteq \cE_{\lambda t}$ and
\begin{align*}
   \big\{E\in \cE_{ct}&\colon \aAd{f+g}{ct} \ge ct\big\}\subseteq  \big\{E\in \cE_{ct}\colon \aAd{f}{ct}+\aAd{f}{ct} \ge t\big\}
   \\ &\subseteq\big\{E\in \cE_{ct}\colon \aAd{f}{ct} \ge \lambda t\}\cup\{E\in \cE_{ct}\colon \aAd{g}{ct} \ge (1-\lambda)t\big\}
    \\ &\subseteq\big\{E\in \cE_{\lambda t}\colon \aAd{f}{\lambda t} \ge \lambda t\big\}\cup \big\{E\in \cE_{(1-\lambda) t}\colon \aAd{g}{(1-\lambda) t} \ge (1-\lambda)t\big\}.
\end{align*}
Hence, $\bmub(f+g,ct)\le \bmub(f,\lambda t)\vee \bmub(g,(1-\lambda) t),$ as claimed.
\end{proof}

\end{document}
